\chapter{Εισαγωγή}
\label{ch:intro}

\section{Αντικείμενο της εργασίας}
\label{sec:intro-scope}

Η εργασία εξετάζει την επιτάχυνση της επιλογής \tl{top-k}, δηλαδή του
εντοπισμού των \tl{k} μεγαλύτερων στοιχείων ενός συνόλου \tl{n} τιμών, σε τρεις
διαφορετικές αρχιτεκτονικές: έναν επεξεργαστή γενικού σκοπού, μια μονάδα
γραφικών και μια μονάδα νευρωνικής επεξεργασίας ενσωματωμένη στο ίδιο πακέτο με
τον επεξεργαστή.

Για κάθε αρχιτεκτονική υλοποιούνται δύο αλγόριθμοι με θεμελιωδώς διαφορετική
συμπεριφορά: η περικομμένη διτονική ταξινόμηση, η οποία εκτελεί σταθερή
ακολουθία συγκρίσεων ανεξάρτητη από τα δεδομένα και η επιλογή με τοπικούς
σωρούς κατά το σχήμα \tl{map-reduce}, η οποία προσαρμόζει την εργασία της στις
τιμές της εισόδου. Και οι έξι υλοποιήσεις μετρώνται στον ίδιο παραμετρικό χώρο,
με κοινή υποδομή, κοινό έλεγχο ορθότητας και κοινή σημασιολογία μέτρησης, ώστε οι
διαφορές που προκύπτουν να αποδίδονται στους αλγορίθμους και στο υλικό και όχι
στον τρόπο μέτρησης.

Το ερώτημα που τίθεται δεν είναι απλώς ποια αρχιτεκτονική είναι ταχύτερη. Είναι
αν και υπό ποιες συνθήκες συμφέρει η ανάθεση του υπολογισμού σε επιταχυντή, όταν
προσμετρηθεί ολόκληρο το κόστος της ανάθεσης και όχι μόνο ο χρόνος του πυρήνα.
Όπως θα φανεί, οι δύο διατυπώσεις του ερωτήματος δίνουν αντίθετες απαντήσεις.

\section{Κίνητρο}
\label{sec:intro-motivation}

Οι μονάδες νευρωνικής επεξεργασίας έπαψαν πρόσφατα να αποτελούν εξοπλισμό
κέντρων δεδομένων. Ενσωματώνονται πλέον στον ίδιο κρύσταλλο με τον επεξεργαστή
φορητών υπολογιστών ευρείας κατανάλωσης, μοιράζονται την ίδια κύρια μνήμη και
είναι διαθέσιμες σε κάθε εφαρμογή που θα μπορούσε να τις αξιοποιήσει. Το ερώτημα
«αξίζει να ανατεθεί αυτός ο υπολογισμός στη μονάδα?» μετατοπίστηκε έτσι από τον
σχεδιασμό υποδομών στην καθημερινή πρακτική του προγραμματιστή.

Η απάντηση δεν είναι προφανής, διότι οι μονάδες αυτές σχεδιάστηκαν για
συγκεκριμένο σκοπό. Το οικοσύστημα λογισμικού τους είναι οργανωμένο γύρω από
τελεστές νευρωνικών δικτύων και η αρχιτεκτονική τους βασίζεται στην παραδοχή
ότι κάθε δεδομένο που μεταφέρεται θα χρησιμοποιηθεί πολλές φορές. Τι συμβαίνει
όταν το φορτίο δεν έχει αυτή τη μορφή παραμένει σε μεγάλο βαθμό ανεξερεύνητο.

Η επιλογή \tl{top-k} αποτελεί κατάλληλο μέσο για να τεθεί το ερώτημα. Είναι
πρόβλημα απλό στη διατύπωση, αλλά διαδεδομένο στην πράξη, καθώς εμφανίζεται στην
αναζήτηση πλησιέστερων γειτόνων, στα συστήματα ανάκτησης και συστάσεων και στη
δειγματοληψία των γλωσσικών μοντέλων. Ταυτόχρονα βρίσκεται στο αντίθετο άκρο του
φάσματος από τα φορτία για τα οποία σχεδιάστηκαν οι επιταχυντές: εκτελεί
ελάχιστες πράξεις ανά στοιχείο και δεν επαναχρησιμοποιεί σχεδόν τίποτα. Αν η
ανάθεση συμφέρει ακόμη και εδώ, θα συμφέρει σχεδόν παντού· αν δεν συμφέρει,
έχουμε μια σαφή ένδειξη για τα όρια της χρησιμότητας των μονάδων αυτών.

Δεύτερο, ανεξάρτητο κίνητρο είναι μεθοδολογικό. Οι συγκρίσεις επιταχυντών
αναφέρουν συνήθως τον χρόνο εκτέλεσης του πυρήνα, αποκρύπτοντας το κόστος
μεταφοράς των δεδομένων από και προς τη συσκευή. Το κόστος αυτό είναι όμως
ακριβώς εκείνο που κρίνει αν η ανάθεση συμφέρει. Η εργασία μετρά συστηματικά και
τα δύο, και δείχνει ότι η επιλογή του ορίου μέτρησης δεν μεταβάλλει απλώς τα
νούμερα, αλλά αντιστρέφει πλήρως την κατάταξη των αρχιτεκτονικών.

\section{Συνεισφορά}
\label{sec:intro-contribution}

Η συνεισφορά της εργασίας είναι τετραπλή.

\paragraph{Υλοποιήσεις.} Κατασκευάστηκαν έξι υλοποιήσεις, δύο αλγόριθμοι επί
τριών αρχιτεκτονικών, με υποστήριξη πέντε τύπων δεδομένων. Η υλοποίηση στον
επεξεργαστή αξιοποιεί τις διανυσματικές επεκτάσεις \tl{AVX-512} και πολλαπλά
νήματα, εκείνη της μονάδας γραφικών χρησιμοποιεί οργάνωση σε πλακίδια με κοινή
μνήμη, ενώ εκείνη της μονάδας νευρωνικής επεξεργασίας κατασκευάστηκε εξ ολοκλήρου
στο επίπεδο του πλαισίου \tl{MLIR-AIE}, καθώς δεν υπάρχει διαθέσιμη βιβλιοθήκη
επιλογής \tl{top-k} για την αρχιτεκτονική αυτή. Η ορθότητα κάθε συνδυασμού
επαληθεύτηκε σε 2700 περιπτώσεις, όλες επιτυχείς.

\paragraph{Υποδομή μέτρησης.} Αναπτύχθηκε υποδομή που μετρά κάθε περίπτωση σε δύο
ρητά διακριτά πεδία, το αλγοριθμικό και το \tl{end-to-end}. Καταγράφει ενέργεια
με κοινή σημασιολογία και στις τρεις αρχιτεκτονικές, καταμετρά τον όγκο των
δεδομένων που διακινούνται και τοποθετεί κάθε εκτέλεση σε οροφή που έχει
μετρηθεί στο ίδιο μηχάνημα. Η υποδομή είναι ανεξάρτητη από το συγκεκριμένο πρόβλημα 
και εφαρμόσιμη σε οποιοδήποτε φορτίο εξετάζεται για ανάθεση σε επιταχυντή.

\paragraph{Ευρήματα.} Η πειραματική αποτίμηση καταλήγει σε τέσσερα συμπεράσματα.
Η επιλογή \tl{top-k} είναι πρόβλημα κίνησης δεδομένων σε κάθε αρχιτεκτονική που
εξετάστηκε και καμία υλοποίηση δεν περιορίζεται από την υπολογιστική ικανότητα.
Η κατάταξη των δύο αλγορίθμων αντιστρέφεται με τον λόγο $n/k$ και όχι με την
αρχιτεκτονική, με διαφορά που φτάνει τις τρεις τάξεις μεγέθους. Η επιλογή του
ορίου μέτρησης αντιστρέφει πλήρως την κατάταξη των αρχιτεκτονικών. Τέλος,
χαμηλή κατανάλωση ισχύος δεν συνεπάγεται χαμηλή κατανάλωση ενέργειας: η μονάδα
νευρωνικής επεξεργασίας έχει τη χαμηλότερη ισχύ από όλες τις υλοποιήσεις και
ταυτόχρονα καταναλώνει πολλαπλάσια ενέργεια (διότι χρειάζεται πολλαπλάσιο χρόνο).

\paragraph{Αναπαραγωγιμότητα.} Ολόκληρη η πειραματική διαδικασία, από την
κατασκευή των εκτελέσιμων έως την παραγωγή των διαγραμμάτων, εκτελείται από
δέσμες εντολών που περιλαμβάνονται στο αποθετήριο. Οι συνθήκες μέτρησης
σταθεροποιούνται ρητά και καταγράφονται μαζί με τα αποτελέσματα. Οι λεπτομέρειες
δίνονται στο παράρτημα~\ref{app:reproducibility}.

\section{Δομή της εργασίας}
\label{sec:intro-outline}

Το κεφάλαιο~\ref{ch:background} ορίζει το πρόβλημα, παρουσιάζει τις οικογένειες
αλγορίθμων που το λύνουν και περιγράφει τα εκτελεστικά μοντέλα των τριών
αρχιτεκτονικών, καταλήγοντας στις υπάρχουσες υλοποιήσεις και στο κενό που η
εργασία καλείται να καλύψει.

Το κεφάλαιο~\ref{ch:methodology} περιγράφει το σύστημα αναφοράς και την
πειραματική διάταξη: τι ακριβώς μετριέται, με ποια σημασιολογία, με ποια μέτρα
σύγκρισης και υπό ποιες συνθήκες. Το κεφάλαιο αυτό είναι προϋπόθεση για την
ανάγνωση των αποτελεσμάτων.

Τα κεφάλαια~\ref{ch:cpu},~\ref{ch:gpu} και~\ref{ch:npu} παρουσιάζουν τις
υλοποιήσεις σε καθεμία από τις τρεις αρχιτεκτονικές. Ακολουθούν κοινή δομή, ώστε
οι αντίστοιχες ενότητες να μπορούν να συγκριθούν άμεσα: σχεδιαστικές επιλογές,
οι δύο αλγόριθμοι, διαχείριση μνήμης και κίνηση δεδομένων, υλοποίηση αναφοράς
και, τέλος, οι περιοριστικοί παράγοντες που εντοπίστηκαν.

Το κεφάλαιο~\ref{ch:evaluation} παρουσιάζει και ερμηνεύει τα αποτελέσματα:
ορθότητα, χρόνο εκτέλεσης, σύγκριση αρχιτεκτονικών, ανάλυση με το μοντέλο
\tl{roofline}, ενεργειακή αποτίμηση και, τέλος, την επίδραση του τύπου δεδομένων
και της κατανομής της εισόδου. Το κεφάλαιο~\ref{ch:conclusions} συνοψίζει όσα
γενικεύονται, δηλώνει τα όρια εντός των οποίων ισχύουν και προτείνει κατευθύνσεις
συνέχειας.

Τα παραρτήματα περιλαμβάνουν κατάλογο ακρωνυμίων, γλωσσάρι των ελληνικών όρων με
τους αγγλικούς αντίστοιχούς τους και οδηγίες αναπαραγωγής των μετρήσεων.
